% =========================
% puzzles/pXX.tex
% Final - The Mirror Lemma Key
% =========================

\PuzzleChapterIntro{The Mirror Lemma Key}{Geometry \textbullet\ Lemma Chain \textbullet\ Unique Answer}{This puzzle connects geometry to graph theory. By translating the 'mirror-true' property into integer parity, you can bound the maximum number of specific triangles using the dual tree of the triangulation.}

\Lore{
On a high hook in the Vault hangs a ring of glass like a crown, its rim cut into \(2026\) facets that catch the lamp in cold, steady flashes.

A pilgrim named Irin is given charcoal and a straightedge and told to cut the ring into triangles, nothing more.

When a triangle settles the glass in just the right way, the crown answers with a clean tone that seems to come from inside the material.

Each tone advances the Mirror Lemma Key by a single click.

The ring does not care how elegant the drawing is---only how many true tones can be forced.
}

\Problem

Let \(\Omega\) be a regular \(2026\)-gon.

A segment joining two vertices of \(\Omega\) (either a side or a diagonal) is called \textbf{mirror-true}
if its endpoints split the boundary of \(\Omega\) into two arcs, each consisting of an \emph{odd} number of sides of \(\Omega\).

Irin draws \(2023\) diagonals of \(\Omega\), no two intersecting in the interior, thereby dissecting \(\Omega\) into \(2024\) triangles.

A triangle in the dissection is called a \textbf{true-tone} triangle if:
\begin{itemize}
  \item it is isosceles, and
  \item exactly two of its sides are mirror-true.
\end{itemize}

\VaultDirective{What is the \textbf{maximum possible} number of true-tone triangles in such a dissection?}


\SolutionPage

Label vertices \(1,2,\dots,2026\) cyclically, \(N=2026\). For \(i,j\), let \(s(i,j)=\min(|i-j|,N-|i-j|)\).
The chord \(ij\) splits the boundary into arcs of lengths \(s(i,j)\) and \(N-s(i,j)\), which have the same parity since \(N\) is even.
So \(ij\) is mirror-true iff \(s(i,j)\) is odd iff \(|i-j|\) is odd, i.e., iff \(i,j\) have opposite parity. In particular, every polygon side is mirror-true.

\smallskip
In any triangle, either all three vertices share a parity (then \(0\) mirror-true sides) or exactly one differs (then exactly the two incident sides are mirror-true). Hence every triangle has \(0\) or \(2\) mirror-true sides.

\smallskip
In a regular \(N\)-gon, chord length depends only on \(s(i,j)\), so equal sides have step sizes of the same parity. Therefore in a true-tone triangle (two mirror-true sides, one non-true side), the equal sides must be the two mirror-true sides; the base is the unique non-true side, hence a diagonal (not a polygon side).

\smallskip
Let \(T\) be the dual graph of the triangulation (a tree on \(2024\) triangles, \(2023\) edges). Let \(k\) be the number of mirror-true diagonals among the drawn diagonals, and delete from \(T\) the \(k\) edges corresponding to those diagonals, obtaining a forest \(F\) with \(k+1\) components.
A triangle with \(0\) mirror-true sides contains no polygon side (all polygon sides are mirror-true), hence all three of its sides are diagonals, all non-true; so it has degree \(3\) in \(F\).
A triangle with \(2\) mirror-true sides has exactly one non-true side, necessarily a diagonal, so has degree \(1\) in \(F\).
Thus every component of \(F\) is a \(\{1,3\}\)-degree tree. If such a tree has \(n\) vertices and \(L\) leaves, then
\(L+3(n-L)=2(n-1)\Rightarrow L=\frac{n+2}{2}\).
Summing over components (\(\sum n=2024\)) gives total leaves
\[
\sum \frac{n_i+2}{2}=\frac{2024+2(k+1)}{2}=1013+k,
\]
so there are \(1013+k\) triangles with \(2\) mirror-true sides.

\smallskip
A mirror-true diagonal \(XY\) cannot border two true-tone triangles. Let \(s=s(X,Y)\), which is odd. In a true-tone triangle, the mirror-true sides are the equal sides, so \(XY\) must be an equal side and the apex is \(X\) or \(Y\). If the apex is \(X\), then the third vertex must be the other vertex (if any) at step size \(s\) from \(X\); if the apex is \(Y\), it must be the other vertex (if any) at step size \(s\) from \(Y\). In both cases that third vertex lies on the same boundary arc between \(X\) and \(Y\), hence on the same side of chord \(XY\). Since the dissection has two triangles adjacent to \(XY\) on opposite sides, at most one is true-tone. Hence at least \(k\) of the \(1013+k\) leaves are not true-tone, so
\[
\#(\text{true-tone})\le (1013+k)-k=1013.
\]

\smallskip
Construction: draw the \(1013\) noncrossing diagonals \((1,3),(3,5),\dots,(2025,1)\), carving off the \(1013\) ear triangles \((1,2,3),(3,4,5),\dots,(2025,2026,1)\).
Each ear is isosceles with exactly two mirror-true sides (the polygon sides). The remaining region uses only odd vertices, so has \(0\) mirror-true sides in every triangle. Thus \(1013\) is achievable.

Therefore the maximum is \(\boxed{1013}\).

\FinalAnswer{1013}

\NextPage